% ==========================================================
% Appendix: Interview Guides (Pilot)
% ==========================================================
\chapter{Pilot interview guides}
\label{app:interview_guides}

This appendix provides the semi-structured interview guides used in the pilot phase. The guides are designed to elicit: (i) perceived challenges across modelling activities (RQ1), (ii) coping strategies and articulated support needs (RQ2), and (iii) socio-technical shaping factors (RQ3). Questions are asked flexibly; probes are used selectively depending on time and participant responses.

\section{Instructor pilot interview guide (20--30 min)}
\label{app:guide_instructor}

\subsection*{0) Opening (30--45s)}
\textbf{Script (say verbatim).}
\begin{quote}
Thanks for meeting with me. I’m studying how modelling is taught and learned in higher education, focusing on challenges, coping strategies/support needs, and how context (tools, course design, domains, feedback) shapes the experience. With your permission, I’d like to record so I don’t miss details. The interview should take about 20--30 minutes. You can skip any question.
\end{quote}

\subsection*{1) Teaching context (3--4 min)}
\textbf{Q1. Teaching background and where modelling fits.}
\begin{itemize}
  \item \emph{Main prompt:} Could you briefly describe your teaching context and where modelling fits into your courses?
  \item \emph{Probes (optional):} courses/level; modelling families (UML/BPMN/ER/iStar/MDE); years taught; class size; team vs individual work.
\end{itemize}

\subsection*{2) Learning goals and course design (6--7 min)}
\textbf{Q2. Intended modelling competence.}
\begin{itemize}
  \item \emph{Main prompt:} When a student finishes your course, what should they be able to do with modelling?
  \item \emph{Probes:} key activities (interpret/create/refine/evaluate); minimum vs excellent; what “good enough” looks like.
\end{itemize}

\textbf{Q3. Scaffolding from problem statement to model (and better model).}
\begin{itemize}
  \item \emph{Main prompt:} How do you scaffold students from the problem statement to a model and a better model?
  \item \emph{Probes:} explicit step-by-step approach; where students get stuck; quality criteria taught explicitly vs via feedback; reflection/comparison tasks.
\end{itemize}

\textbf{Q4. Assessment focus across modelling activities.}
\begin{itemize}
  \item \emph{Main prompt:} How do you assess modelling competence in your course?
  \item \emph{Probes:} final artefact vs process/reasoning; which activities assessed; quality dimensions graded (completeness/consistency/readability/fit-for-purpose).
\end{itemize}

\subsection*{3) Observed student challenges (6--8 min)}
\textbf{Q5. Biggest recurring challenges.}
\begin{itemize}
  \item \emph{Main prompt:} What are the most common challenges students face in modelling?
  \item \emph{Probes:} conceptual (semantics) vs procedural (how to build/refine); which activity is hardest; differences by notation; beginner vs advanced differences.
\end{itemize}

\textbf{Q6. Concrete incident walkthrough (high value).}
\begin{itemize}
  \item \emph{Main prompt:} Can you recall a specific moment when students struggled with modelling? What was happening?
  \item \emph{Probes:} what exactly stuck; how you noticed; conditions (tool, time pressure, group work, unclear requirements, domain unfamiliarity); redesign ideas.
\end{itemize}

% Optional: keep or drop depending on time
\subsection*{Optional mini-elicitation (2--3 min)}
\textbf{Prompt.}
\begin{itemize}
  \item Could you briefly describe one modelling assignment/exam task you use, and what a strong vs weak submission looks like?
  \item \emph{Probes:} which activities it tests; diagnostic mistakes; what feedback/support would prevent them.
\end{itemize}

\subsection*{4) Coping strategies and support needs (6--7 min)}
\textbf{Q7. What you do when students struggle.}
\begin{itemize}
  \item \emph{Main prompt:} When students struggle, what do you typically do to help, and what seems to work best?
  \item \emph{Probes:} worked examples/templates/checklists/peer review; tailoring by difficulty type; what students do effectively/ineffectively.
\end{itemize}

\textbf{Q8. Tool-related supports and constraints.}
\begin{itemize}
  \item \emph{Main prompt:} What role do tools play in these challenges, and what tool support would be most helpful?
  \item \emph{Probes:} syntax checking, hints, layout/views, versioning/collab; tool friction (setup/usability/cognitive load); tool changes due to student experience.
\end{itemize}

\textbf{Q9. Feedback and iteration.}
\begin{itemize}
  \item \emph{Main prompt:} How do you provide feedback on models, and how much opportunity do students have to revise?
  \item \emph{Probes:} one-shot vs iterative; explicit quality criteria; approaches that improved students’ evaluate/refine ability.
\end{itemize}

\subsection*{5) Socio-technical shaping factors (3--4 min)}
\textbf{Q10. Domains and authenticity.}
\begin{itemize}
  \item \emph{Main prompt:} How do you choose assignment domains/examples, and how do those choices affect engagement or performance?
  \item \emph{Probes:} familiar vs unfamiliar domains; allowing domain choice; demotivating/excluding domain choices.
\end{itemize}

\subsection*{6) Closing (2--3 min)}
\textbf{Q11. Ideal supports / recommendations.}
\begin{itemize}
  \item \emph{Main prompt:} If you could change one or two things to improve modelling learning in your context, what would you change?
  \item \emph{Probes:} course design/scaffolding/assessment; feedback; tool features; “small change, big impact”.
\end{itemize}

\textbf{Q12. Anything missing.}
\begin{itemize}
  \item \emph{Main prompt:} Is there anything important I didn’t ask but should have?
\end{itemize}

\bigskip
\noindent\textbf{Close (say verbatim).}
\begin{quote}
Thank you — this was very helpful. If you’re interested, I can share a short summary of findings later, and I may follow up with one or two clarifying questions.
\end{quote}

\section{Student pilot interview guide (20--30 min)}
\label{app:guide_student}

\subsection*{0) Opening (1--2 min)}
\textbf{Script (say verbatim).}
\begin{quote}
Hi! Thanks for taking the time. I’m doing a thesis on challenges in learning and doing software/conceptual modelling in higher education. This will take about 20--30 minutes. There are no right or wrong answers—I’m interested in your experience. With your permission, I’d like to audio-record so I don’t miss details. The recording will be anonymised and used only for research. You can skip any question, and you can stop the interview at any point. Is it okay if I record?
\end{quote}

\subsection*{1) Background and context (3--4 min)}
\textbf{Q1. Studies and modelling exposure.}
\begin{itemize}
  \item \emph{Main prompt:} Can you tell me a bit about your studies and where modelling fits in (programme, year, courses)?
  \item \emph{Probes:} which courses and purposes; notations used; modelling outside courses (work/internship/personal projects).
\end{itemize}

\subsection*{2) Experience with modelling (6--7 min)}
\textbf{Q2. Personal meaning of modelling.}
\begin{itemize}
  \item \emph{Main prompt:} When you hear “software modelling” or “conceptual modelling”, what comes to your mind?
  \item \emph{Probes:} first example/image; useful work vs extra documentation; 1--2 sentence explanation to a peer.
\end{itemize}

\textbf{Q3. Useful moment story.}
\begin{itemize}
  \item \emph{Main prompt:} Can you remember a moment when modelling felt particularly useful?
  \item \emph{Probes:} context/task; what it helped (understand/decide/communicate); effect on confidence/motivation.
\end{itemize}

\textbf{Q4. Frustrating moment story.}
\begin{itemize}
  \item \emph{Main prompt:} Can you remember a moment when modelling felt confusing, pointless, or frustrating?
  \item \emph{Probes:} cause (notation/tool/task/feedback/group/time); coping response; what would have helped.
\end{itemize}

\subsection*{3) Learning and challenges across modelling activities (7--8 min)}
\textbf{Q5. Learning a new notation/framework.}
\begin{itemize}
  \item \emph{Main prompt:} When you learn a new notation/framework, what’s usually the hardest part at the beginning?
  \item \emph{Probes:} symbols vs semantics vs how to start; specific notation; strategies used to learn.
\end{itemize}

\textbf{Q6. Creating a model from a text description.}
\begin{itemize}
  \item \emph{Main prompt:} Imagine you get a text description and have to build a model. What steps do you take?
  \item \emph{Probes:} starting point; structured strategy vs trial-and-error; what to include/omit; whether taught a step-by-step approach.
\end{itemize}

\textbf{Q7. Deciding whether a model is “good enough”.}
\begin{itemize}
  \item \emph{Main prompt:} Once you have a model, how do you decide whether it’s good enough or needs changes?
  \item \emph{Probes:} criteria (completeness/consistency/readability/fit); quality rules/checklists; comparing with peers.
\end{itemize}

\subsection*{4) Socio-technical factors (6--8 min)}
\textbf{Q8. Tools.}
\begin{itemize}
  \item \emph{Main prompt:} What tools have you used for modelling, and how has that been?
  \item \emph{Probes:} liked/disliked tools and why; tool makes modelling harder/easier; whether tools helped catch mistakes.
\end{itemize}

\textbf{Q9. Assignments, domains, authenticity, and inclusion.}
\begin{itemize}
  \item \emph{Main prompt:} What kinds of domains/examples did you work with in modelling assignments?
  \item \emph{Probes:} motivating vs alien domains; domain choice; gentle inclusion probe (if appropriate).
\end{itemize}

\textbf{Q10. Feedback and iteration.}
\begin{itemize}
  \item \emph{Main prompt:} How did feedback on modelling work in your courses?
  \item \emph{Probes:} type/timing of feedback; what would have helped most; effect of group work on modelling.
\end{itemize}

\subsection*{5) Closing (3--4 min)}
\textbf{Q11. “Good modeller” competence.}
\begin{itemize}
  \item \emph{Main prompt:} What should a good modeller be able to do?
  \item \emph{Probes:} habits/abilities; soft skills vs technical; whether courses helped develop them.
\end{itemize}

\textbf{Q12. Improvements and support needs.}
\begin{itemize}
  \item \emph{Main prompt:} If you could change one or two things in how modelling is taught, what would you change?
\end{itemize}

\textbf{Final open question.}
\begin{itemize}
  \item Is there anything important about learning or using modelling that we didn’t cover?
\end{itemize}

\section{Mapping of interview questions to research questions}
\label{app:rq_mapping}

Table~\ref{tab:rq_mapping_guides} summarises how the main interview prompts operationalise the analytic focus of each research question. Probes are used flexibly to adapt to participant experience and time constraints.

\begin{table}[htbp]
\centering
\small
\setlength{\tabcolsep}{6pt}
\renewcommand{\arraystretch}{1.15}
\caption{Mapping of pilot interview guide questions to research questions.}
\label{tab:rq_mapping_guides}
\begin{tabularx}{\textwidth}{l X c c c}
\toprule
\textbf{Guide} & \textbf{Question (main prompt)} & \textbf{RQ1} & \textbf{RQ2} & \textbf{RQ3} \\
\midrule
Student & Q1 Background and modelling exposure &  &  & \checkmark \\
Student & Q2 Meaning/perception of modelling & \checkmark &  & \checkmark \\
Student & Q3 Useful moment (incident) & \checkmark & \checkmark & \checkmark \\
Student & Q4 Frustrating moment (incident) & \checkmark & \checkmark & \checkmark \\
Student & Q5 Learning a new notation/framework & \checkmark & \checkmark & \checkmark \\
Student & Q6 Creating a model from text (procedure) & \checkmark & \checkmark &  \\
Student & Q7 Judging quality / “good enough” & \checkmark & \checkmark &  \\
Student & Q8 Tools & \checkmark & \checkmark & \checkmark \\
Student & Q9 Domains/authenticity/inclusion &  & \checkmark & \checkmark \\
Student & Q10 Feedback and iteration & \checkmark & \checkmark & \checkmark \\
Student & Q11 “Good modeller” competence & \checkmark &  & \checkmark \\
Student & Q12 Improvements / supports &  & \checkmark & \checkmark \\
\midrule
Instructor & Q1 Teaching context and where modelling fits &  &  & \checkmark \\
Instructor & Q2 Intended modelling competence & \checkmark &  & \checkmark \\
Instructor & Q3 Scaffolding from problem to model & \checkmark & \checkmark & \checkmark \\
Instructor & Q4 Assessment focus and criteria & \checkmark & \checkmark & \checkmark \\
Instructor & Q5 Recurring student challenges & \checkmark &  & \checkmark \\
Instructor & Q6 Concrete incident walkthrough & \checkmark & \checkmark & \checkmark \\
Instructor & Q7 Instructor interventions and what works &  & \checkmark & \checkmark \\
Instructor & Q8 Tool supports and constraints &  & \checkmark & \checkmark \\
Instructor & Q9 Feedback and revision opportunities & \checkmark & \checkmark & \checkmark \\
Instructor & Q10 Domains/authenticity effects &  & \checkmark & \checkmark \\
Instructor & Q11 Ideal supports / recommendations &  & \checkmark & \checkmark \\
Instructor & Q12 Anything missing (catch-all) & \checkmark & \checkmark & \checkmark \\
\bottomrule
\end{tabularx}
\end{table}
